% !TEX root = task-1-15-4.tex
% Editable task bank. Append new BankTask blocks at the end of the document.
\documentclass[12pt,a4paper]{article}
\usepackage{physics-textbook}
\usepackage{graphicx,needspace}
\newif\ifShowSolutions
\ifdefined\TutorEdition
  \ShowSolutionstrue
\else
  \ShowSolutionsfalse % Change to \ShowSolutionstrue to display solutions.
\fi
\newcounter{banktask}
\newcommand{\BankTask}[3]{%
  \refstepcounter{banktask}%
  \Needspace{35mm}\par\medskip\noindent\begin{minipage}{\linewidth}%
  \noindent\textbf{\thebanktask.}\quad #2\par
  \ifShowSolutions
    #3\par
    \ifdefined\TutorEdition
      \ifcsname TutorSource#1\endcsname
        \smallskip{\small\csname TutorSource#1\endcsname\par}%
      \fi
    \fi
  \fi
  \end{minipage}\par\medskip
}
\begin{document}
\BookSubsection{1.15.4}{Кинетическая энергия}


% Source sdamgia-527: https://phys-ege.sdamgia.ru/problem?id=527
\BankTask{sdamgia-527}{
Тележка движется со скоростью $3\,\text{м/с}$. Её кинетическая энергия равна $27\,\text{Дж}$. Какова масса тележки?
}{
\smallskip\noindent\textbf{Решение}\par
Кинетическая энергия тележки связана с массой тележки и скоростью её движения соотношением: $E_{\text{кин}}=\frac{mv^2}{2}$. Таким образом, масса тележки равна:
\[m=\frac{2E_{\text{кин}}}{v^2}=\frac{2\cdot27\,\text{Дж}}{(3\,\text{м/с})^2}=6\,\text{кг}.\]
Ответ: $6\,\text{кг}$.
}

% Source sdamgia-536: https://phys-ege.sdamgia.ru/problem?id=536
\BankTask{sdamgia-536}{
Хоккейная шайба массой $160\,\text{г}$ летит со скоростью $10\,\text{м/с}$ без вращения. Какова её кинетическая энергия?
}{
\smallskip\noindent\textbf{Решение}\par
Кинетическая энергия шайбы связана с массой шайбы и скоростью её движения соотношением:
\[E_{\text{кин}}=\frac{mv^2}{2}=\frac{0{,}16\,\text{кг}\cdot(10\,\text{м/с})^2}{2}=8\,\text{Дж}.\]
Ответ: $8\,\text{Дж}$.
}

% Source sdamgia-736: https://phys-ege.sdamgia.ru/problem?id=736
\BankTask{sdamgia-736}{
Скорость брошенного мяча непосредственно перед ударом об абсолютно гладкую стену была вдвое больше его скорости сразу после удара. Какое количество теплоты выделилось при ударе, если перед ударом кинетическая энергия мяча была равна $20\,\text{Дж}$?
}{
\smallskip\noindent\textbf{Решение}\par
Кинетическая энергия пропорциональна квадрату скорости. Поскольку скорость после удара уменьшилась в $2$ раза, кинетическая энергия мяча уменьшилась в $4$ раза и стала равна $5\,\text{Дж}$. По закону сохранения количество теплоты, выделившееся после удара, равняется убыли кинетической энергии:
\[Q=|\Delta E_{\text{кин}}|=E_{\text{кин}}^{(1)}-E_{\text{кин}}^{(2)}=20\,\text{Дж}-5\,\text{Дж}=15\,\text{Дж}.\]
Ответ: $15$.
\par Примечание. Для приведённого здесь решения задачи существенно, что стенка, о которую ударяется мяч, абсолютно гладкая. Действительно, в противном случае на мяч во время удара может дейстовать еще и сила трения со стороны стенки, которая приведёт к тому, что мяч после удара будет вращаться, поэтому кинетическая энергия шара после удара будет состоять из энергии поступательного и вращательного движений. Следовательно, судить о величине выделившегося тепла будет невозможно.
}

% Source sdamgia-3787: https://phys-ege.sdamgia.ru/problem?id=3787
\BankTask{sdamgia-3787}{
Изменение скорости тела массой $2\,\text{кг}$, движущегося по оси $x$, описывается формулой $v_x=v_{0x}+a_xt$, где $v_{0x}=8\,\text{м/с}$, $a_x=-2\,\text{м/с}^2$, $t$ --- время в секундах. Какова кинетическая энергия тела через $3\,\text{с}$ после начала отсчёта времени?
}{
\smallskip\noindent\textbf{Решение}\par
Кинетическая энергия тела связана с его массой и скоростью соотношением: $E_{\text{кин}}=\frac{mv^2}{2}$. Определим скорость тела через $3\,\text{с}$ после начала отсчёта времени:
\[v_x(3\,\text{с})=8\,\text{м/с}-2\,\text{м/с}^2\cdot3\,\text{с}=2\,\text{м/с}.\]
Таким образом, искомая кинетическая энергия равна:
\[E_{\text{кин}}=\frac{2\,\text{кг}\cdot(2\,\text{м/с})^2}{2}=4\,\text{Дж}.\]
Ответ: $4\,\text{Дж}$.
}

% Source sdamgia-4763: https://phys-ege.sdamgia.ru/problem?id=4763
\BankTask{sdamgia-4763}{
Самосвал массой $m_0$ при движении на пути к карьеру имеет кинетическую энергию $2{,}5\cdot10^5\,\text{Дж}$. Какова его кинетическая энергия после загрузки, если он двигался с прежней скоростью, а масса его увеличилась в $2$ раза?
}{
\smallskip\noindent\textbf{Решение}\par
Кинетическая энергия пропорциональна массе тела: $E=\frac{mv^2}{2}$. Таким образом, после загрузки при движении с такой же скоростью самосвал будет иметь вдвое большую кинетическую энергию:
\[2\cdot2{,}5\cdot10^5\,\text{Дж}=5\cdot10^5\,\text{Дж}=500\,\text{кДж}.\]
Ответ: $500\,\text{кДж}$.
}

% Source sdamgia-9016: https://phys-ege.sdamgia.ru/problem?id=9016
\BankTask{sdamgia-9016}{
Два тела движутся с одинаковой скоростью. Кинетическая энергия первого тела в $4$ раза меньше кинетической энергии второго тела. Определите отношение $\frac{m_1}{m_2}$ масс тел.
}{
\smallskip\noindent\textbf{Решение}\par
Исходя из условия, запишем соотношение для кинетических энергий тел:
\[4\cdot\frac{m_1v^2}{2}=\frac{m_2v^2}{2}\quad\Longleftrightarrow\quad\frac{m_1}{m_2}=0{,}25.\]
Ответ: $0{,}25$.
}

% Source sdamgia-3582: https://phys-ege.sdamgia.ru/problem?id=3582
\BankTask{sdamgia-3582}{
Небольшое тело массой $0{,}2\,\text{кг}$ бросили вертикально вверх. На рисунке показан график зависимости кинетической энергии $E_{\text{к}}$ тела от времени $t$ в течение полёта. Чему равна максимальная скорость тела в первые четыре секунды полёта?\par\begin{center}
\begin{tikzpicture}[task graph,x=1cm,y=.1cm]
\draw[task grid,xstep=.5,ystep=5] (0,0) grid (5,50);
\draw[task axis] (-.2,0)--(5.15,0) node[below] {$t,\,\text{с}$};
\draw[task axis] (0,-2)--(0,51) node[left] {$E_{\text{к}},\,\text{Дж}$};
\node[below left] at (0,0) {$0$};
\foreach \x in {1,2,3,4} {\fill (\x,0) circle[x radius=1.2pt,y radius=1.2pt]; \node[below] at (\x,0) {$\x$};}
\foreach \y in {10,20,30,40} {\fill (0,\y) circle[x radius=1.2pt,y radius=1.2pt]; \node[left] at (0,\y) {$\y$};}
\draw[task curve,domain=0:4,samples=81] plot (\x,{10*(\x-2)^2});
\end{tikzpicture}\end{center}
}{
\smallskip\noindent\textbf{Решение}\par
Кинетическая энергия тела связана с его скоростью соотношением $E_{\text{к}}=\frac{mv^2}{2}$. Кинетическая энергия максимальна тогда же, когда максимальна и скорость. Из графика видно, что максимальная кинетическая энергия равна $40\,\text{Дж}$. Следовательно, максимальная скорость равна
\[v_{\max}=\sqrt{\frac{2E_{\text{к}\max}}{m}}=\sqrt{\frac{2\cdot40\,\text{Дж}}{0{,}2\,\text{кг}}}=20\,\text{м/с}.\]
Ответ: $20\,\text{м/с}$.
}

% Source sdamgia-9517: https://phys-ege.sdamgia.ru/problem?id=9517
\BankTask{sdamgia-9517}{
Камень массой $40\,\text{г}$ брошен под углом $60^\circ$ к горизонту. Начальная кинетическая энергия камня равна $2\,\text{Дж}$. Чему равен модуль импульса камня в верхней точке траектории его движения? Сопротивление воздуха пренебрежимо мало.
}{
\smallskip\noindent\textbf{Решение}\par
При отсутствии сопротивления воздуха камень движется только под действием силы тяжести с ускорением свободного падения. Его проекция на ось $Ox$ равна $g_x=0$, следовательно, движение камня в горизонтальном направлении равномерное: $v_x=v_{0x}$. Проекции скорости в верхней точке траектории равны $v_{0x}=v_0\cos\alpha$, $v_y=0$, так как в верхней точке траектории скорость камня направлена горизонтально. Таким образом, скорость камня в верхней точке траектории равна $v=v_x=v_0\cos\alpha$. Кинетическая энергия камня в момент броска $E_{\text{к}}=\frac{mv_0^2}{2}$, откуда его начальная скорость $v_0=\sqrt{\frac{2E_{\text{к}}}{m}}$. Тогда импульс камня в верхней точке траектории равен:
\[\begin{aligned}p=mv=mv_0\cos\alpha&=\sqrt{2E_{\text{к}}m}\cos\alpha\\&=\sqrt{2\cdot2\cdot0{,}04}\cdot\cos60^\circ=0{,}2\,\text{кг}\cdot\text{м/с}.\end{aligned}\]
Ответ: $0{,}2\,\text{кг}\cdot\text{м/с}$.
}

% Source sdamgia-3413: https://phys-ege.sdamgia.ru/problem?id=3413
\BankTask{sdamgia-3413}{
Амплитуда малых свободных колебаний пружинного маятника равна $4\,\text{см}$, масса груза --- $400\,\text{г}$, жёсткость пружины --- $40\,\text{Н/м}$. Чему равна максимальная скорость колеблющегося груза?
}{
\smallskip\noindent\textbf{Решение}\par
В силу закона сохранения энергии максимальная кинетическая и максимальная потенциальная энергия пружинного маятника равны:
\[\frac{mv_{\max}^2}{2}=\frac{kx_{\max}^2}{2}.\]
\[v_{\max}=\sqrt{\frac{k}{m}}x_{\max}=\sqrt{\frac{40}{0{,}4}}\cdot0{,}04=0{,}4\,\text{м/с}.\]
Ответ: $0{,}4\,\text{м/с}$.
\par Приведём другое решение. При малых колебаниях максимальная скорость связана с амплитудой колебаний и круговой частотой соотношением $v_{\text{м}}=\omega x_{\text{м}}$. Так как циклическая частота пружинного маятника равна $\omega=\sqrt{\frac{k}{m}}$, получаем, что максимальная скорость равна
\[v_{\text{м}}=\sqrt{\frac{k}{m}}x_{\text{м}}=\sqrt{\frac{40}{0{,}4}}\cdot0{,}04=0{,}4\,\text{м/с}.\]
}

% Source sdamgia-51377: https://phys-ege.sdamgia.ru/problem?id=51377
\BankTask{sdamgia-51377}{
Мальчик бросил маленький мячик под углом к горизонту с горизонтальной поверхности земли. Через время $t=4\,\text{с}$ мячик упал на землю на расстоянии $s=20\,\text{м}$ от места броска. Чему был равен модуль скорости мячика в тот момент, когда его кинетическая энергия была минимальной? Сопротивлением воздуха пренебречь.
}{
\smallskip\noindent\textbf{Решение}\par
В течение полёта горизонтальная проекция $v_x$ скорости мячика не изменялась. В верхней точке траектории вертикальная проекция $v_y$ скорости мячика была равна нулю. Скорость мячика в любой точке траектории $v=\sqrt{v_x^2+v_y^2}$. Поэтому кинетическая энергия мячика $E_{\text{к}}=\frac{mv^2}{2}$ была минимальной, когда он находился в верхней точке траектории. В этот момент скорость мячика была направлена горизонтально и равна по модулю $v_x$. По горизонтали мячик за время $t$ пролетел расстояние $s=v_xt$. Отсюда искомый модуль скорости равен
\[v=v_x=\frac{s}{t}=\frac{20}{4}=5\,\text{м/с}.\]
Ответ: $5\,\text{м/с}$.
}


% Source cF941A: https://ege.fipi.ru/bank/questions.php?search=1&proj=BA1F39653304A5B041B656915DC36B38&qid=cF941A
\BankTask{cF941A}{
При подаче скорость волейбольного мяча составила $18\,\text{м}/\text{с}$. Масса мяча $280\,\text{г}$. Какова кинетическая энергия мяча? Выразите ответ в джоулях.
}{
\smallskip\noindent\textbf{Решение от GPT}\par
$m=280\,\text{г}=0{,}28\,\text{кг}$. Кинетическая энергия $E_{\text{к}}=mv^2/2=0{,}28\cdot18^2/2=45{,}36\,\text{Дж}$. Ответ: $45{,}36\,\text{Дж}$.
}

% Source 7998DB: https://ege.fipi.ru/bank/questions.php?search=1&proj=BA1F39653304A5B041B656915DC36B38&qid=7998DB
\BankTask{7998DB}{
Два тела движутся с одинаковой скоростью. Кинетическая энергия первого тела в $4$ раза меньше кинетической энергии второго тела. Определите отношение $\frac{m_1}{m_2}$ масс тел.
}{
\smallskip\noindent\textbf{Решение от GPT}\par
При одинаковых скоростях $E_{\text{к}1}/E_{\text{к}2}=m_1/m_2=1/4=0{,}25$.
}

% Source 72e2A7: https://ege.fipi.ru/bank/questions.php?search=1&proj=BA1F39653304A5B041B656915DC36B38&qid=72e2A7
\BankTask{72e2A7}{
На рисунке приведены графики зависимости координаты $x$ от времени $t$ для двух точечных тел $1$ и $2$, движущихся по прямой, вдоль которой направлена ось $Ox$. Масса каждого тела равна $0{,}5\,\text{кг}$.
\begin{center}\begin{tikzpicture}[x=.25cm,y=.5cm,>=stealth]
\draw[thin,black!65,xstep=2,ystep=1] (0,0) grid (20,8);
\draw[->] (0,0)--(22,0) node[right] {$t,\,\text{с}$};
\draw[->] (0,0)--(0,9) node[above] {$x,\,\text{м}$};
\foreach \x in {0,4,8,12,16,20} \node[below] at (\x,0) {$\x$};
\foreach \y in {2,4,6} \node[left] at (0,\y) {$\y$};
\draw[very thick] (0,.5)--(20,5.5);\node at (12,3.9) {$1$};
\draw[very thick,domain=0:19.2,samples=60] plot (\x,{1+(\x-10)^2/16});\node at (17,1.7) {$2$};
\end{tikzpicture}\end{center}
\begin{enumerate}
\item Проекция скорости тела $2$ на ось $Ox$ в момент времени $4\,\text{с}$ отрицательна.
\item Тело $1$ движется со скоростью $3\,\text{м}/\text{с}$.
\item Проекция ускорения тела $2$ на ось $Ox$ отрицательна.
\item Кинетическая энергия тела $1$ в момент времени $10\,\text{с}$ равна $2{,}25\,\text{Дж}$.
\item Расстояние между точками встречи тел равно $3\,\text{м}$.
\end{enumerate}
}{
\smallskip\noindent\textbf{Решение от GPT}\par
При $t=4\,\text{с}$ график тела $2$ убывает, значит $v_{2x}<0$. Скорость тела $1$ равна наклону прямой: $v_1=(5-2)/(18-6)=0{,}25\,\text{м}/\text{с}$, а его энергия $E_{\text{к}1}=0{,}5\cdot0{,}25^2/2\ne2{,}25\,\text{Дж}$. Парабола вогнута вверх, поэтому $a_{2x}>0$. Координаты встреч $x_a=2\,\text{м}$ и $x_b=5\,\text{м}$ отличаются на $3\,\text{м}$. Ответ: $15$.
}

% Source 1D0119: https://ege.fipi.ru/bank/questions.php?search=1&proj=BA1F39653304A5B041B656915DC36B38&qid=1D0119
\BankTask{1D0119}{
Шайба массой $m$, скользящая по гладкой горизонтальной поверхности со скоростью $v$, абсолютно неупруго сталкивается с покоящейся шайбой массой $M$.
\TaskMatchingColumns[.53\linewidth]{Физические величины}{
А) кинетическая энергия покоившейся шайбы после столкновения\par
Б) импульс налетающей шайбы после столкновения
}{Формулы}{
1) $\frac{m^2v}{m+M}$\par
2) $mv$\par
3) $\frac{m^2Mv^2}{2(m+M)^2}$\par
4) $\frac{Mv^2}{2}$\par
}\begin{center}\begin{tabular}{|c|c|}\hline А & Б \\\hline\rule{0pt}{1.2em}\hspace{1cm} & \hspace{1cm} \\\hline\end{tabular}\end{center}
}{
\smallskip\noindent\textbf{Решение от GPT}\par
Из сохранения импульса общая скорость $u=mv/(m+M)$. Энергия второй шайбы $Mu^2/2=m^2Mv^2/[2(m+M)^2]$ (формула $3$); импульс первой $mu=m^2v/(m+M)$ (формула $1$). Ответ: $31$.
}

% Source 010677: https://ege.fipi.ru/bank/questions.php?search=1&proj=BA1F39653304A5B041B656915DC36B38&qid=010677
\BankTask{010677}{
Камень подбросили вверх. Как меняются по мере подъёма ускорение камня
и его кинетическая энергия? Сопротивление воздуха не учитывать.
\begin{enumerate}
\item увеличивается
\item уменьшается
\item не меняется
\end{enumerate}
\TaskChangeTable{Ускорение камня}{Кинетическая энергия камня}
}{
\smallskip\noindent\textbf{Решение от GPT}\par
Ускорение свободного падения $\vec a=\vec g$ постоянно. При подъёме скорость уменьшается, поэтому $E_{\text{к}}=mv^2/2$ уменьшается. Ответ: $32$.
}

% Source 52049: https://phys-ege.sdamgia.ru/problem?id=52049
\BankTask{52049}{
Небольшой груз массой m  =  $200\,\text{г}$ совершает незатухающие гармонические колебания на невесомой нерастяжимой нити длиной l  =  $2\,\text{м}$. Закон движения этого груза имеет вид: $x ( t ) = 0{,}2 \cos ( \omega t ) ,$ где величина смещения x от положения равновесия измеряется в метрах. Определите максимальную кинетическую энергию груза в процессе колебаний.
}{
\smallskip\noindent\textbf{Решение}\par
Максимальная кинетическая энергия груза равна: 

\[E_k = \frac{m v _{\max} ^2}{2} , \qquad ( 1 )\]

где максимальная скорость груза в процессе колебаний 

\[v _{\max} = \omega A, \qquad ( 2 )\]

а A  — амплитуда колебаний маятника, величину которой можно найти из закона движения: $A = 0{,}2 \,\text{м}.$
Частота \ensuremath{\omega} и период T колебаний маятника связаны соотношением: 

\[\omega = \frac{2 \pi}{T} . \qquad ( 3 )\]

Период колебаний математического маятника можно найти по формуле: 

\[T = 2 \pi \sqrt{\frac{l}{g}} . \qquad ( 4 )\]

Тогда из уравнений ($1$)–($4$) получим: 

\[E_k = \frac{m g A ^2}{2 l} = \frac{0{,}2 \cdot 10 \cdot 0{,}2 ^2}{2 \cdot 2} = 0{,}02 \text \,\text{Дж}.\]

Ответ: $0{,}02\,\text{Дж}$.
}

% Source 2942: https://phys-ege.sdamgia.ru/problem?id=2942
\BankTask{2942}{
Снаряд массой $4\,\text{кг}$, летящий со скоростью $400\,\text{м/с}$, разрывается на две равные части, одна из которых летит в направлении движения снаряда, а другая — в противоположную сторону. В момент разрыва суммарная кинетическая энергия осколков увеличилась на величину $\Delta E=0{,}5\,\text{МДж}$. Определите скорость осколка, летящего по направлению движения снаряда.
}{
\smallskip\noindent\textbf{Решение}\par
\textbf{Обоснование.}

Задачу будем решать в инерциальной системе отсчета, связанной с Землей. Все тела будем считать материальными точками. Трением снаряда и осколков о воздух пренебрежем.

Для описания разрыва снаряда использован закон сохранения импульса системы тел. Он выполняется в инерциальной системе отсчета, если сумма внешних сил, приложенных к телам системы, равна нулю. В данном случае из-за отсутствия сопротивления воздуха внешней силой является только сила тяжести $m\vec g$, которая не равна нулю. Но этим можно пренебречь, считая время разрыва гранаты малым. За малое время разрыва импульс каждого из осколков меняется на конечную величину за счет больших внутренних сил, разрывающих снаряд при взрыве. По сравнению с этими большими силами конечная сила тяжести пренебрежимо мала. Так как время разрыва гранаты считаем малым, то можно пренебречь и изменением потенциальной энергии гранаты и его осколков в поле тяжести в процессе разрыва. В инерциальной системе отсчета выполняется законы сохранения импульса и энергии с учетом энергии разрыва.

Перейдем к решению.

По закону сохранения импульса тел: $2m\vec v_0=m\vec v_1+m\vec v_2$, где $m$ — масса одного осколка. По закону сохранения энергии для незамкнутой системы: $\Delta E=E_{\text{к}1}+E_{\text{к}2}-E_{\text{к}0}$, где кинетическая энергия осколков после разрыва $E_{\text{к}1}=\frac{mv_1^2}{2}$ и $E_{\text{к}2}=\frac{mv_2^2}{2}$, кинетическая энергия снаряда до взаимодействия $E_{\text{к}0}=\frac{2mv_0^2}{2}$. Проецируя векторную запись на горизонтальную ось, объединяя уравнения, получаем систему уравнений для решения задачи:
\[\begin{cases}2mv_0=mv_1-mv_2,\\\dfrac{2mv_0^2}{2}+\Delta E=\dfrac{mv_1^2}{2}+\dfrac{mv_2^2}{2}.\end{cases}\]
Выразим $v_2$ из первого уравнения: $v_2=v_1-2v_0$ и подставим во второе уравнение. Получим:
\[v_1^2-2v_0v_1+v_0^2-\frac{\Delta E}{m}=0.\]
Отсюда следует: $\Delta E=m(v_1-v_0)^2$. Скорость осколка:
\[v_1=v_0+\sqrt{\frac{\Delta E}{m}}=400+\sqrt{\frac{0{,}5\cdot10^6}{2}}=900\,\text{м/с}.\]
Ответ: $v_1=900\,\text{м/с}$.
}

\end{document}
